\usepackage[a4paper,
        margin=2.5cm,
        top=2.5cm,
        bottom=2.5cm
    ]{geometry}
\usepackage[
    backend=biber,
    style=numeric-comp,
    sorting=none,
    natbib=true
]{biblatex}

\iftoggle{useheaders}{
    \pagestyle{fancy}
    \fancyhf{}
    \fancyhead[R]{\thepage}
    \fancyhead[L]{\slshape\nouppercase{\leftmark}} % Itálico (slanted)
    \renewcommand{\headrulewidth}{0.4pt} % Linha charmosa separando o header
}{
    \pagestyle{plain} 
}

\usepackage{mathpazo}

%> Formatação dos capítulos enumerados
\titleformat{\chapter}[display]
    {\normalfont\bfseries}
    {\filleft\LARGE\chaptertitlename\ \thechapter} % Alinhamento à direita para "Capítulo X"
    {0.5cm} % Espaço entre o número e a linha
    {\titlerule\vspace{0.8cm}\filright\Huge} % Título do capítulo alinhado à esquerda
    [\vspace{2cm}]
  
\titleformat{name=\chapter,numberless}[display]
  {\normalfont\bfseries}{}{0pt}
  {\filright\Huge} % AQUI REMOVEMOS O \titlerule
  [\vspace{2cm}]

\titlespacing*{\chapter}{0pt}{0pt}{0pt}

%> Configuração das epígrafes dos capítulos
\renewcommand{\epigraph}[2]{
  \begin{flushright}
  \begin{minipage}{.3\textwidth}
    \hspace{0.3cm}{#1}\vspace{0.3cm}\\
    {\itshape\small#2}
  \end{minipage}\hspace*{2cm}
\end{flushright}
\vspace{1.5cm}
}